% !TEX root = ../algorithms.tex

\section{Приближённые алгоритмы. Задача о вершинном покрытии минимальной мощности, приближённый полиномиальный алгоритм для её решения. Задача о вершинном покрытии минимальной стоимости, её приближённое решение при помощи лин.~прог.}

Многие важные задачи оптимизации являются NP-полными: для них неизвестно полиномиальных алгоритмов, а нахождение точного ответа, по всей видимости, требует экспоненциального времени. На практике такие задачи обычно решаются либо эвристиками (без гарантий качества), либо \emph{приближёнными алгоритмами} --- полиномиальными методами, гарантирующими ответ, близкий к оптимальному.

\begin{Definition}{Приближённый алгоритм ($\rho$-приближение)}{}
	Пусть дана задача оптимизации $f\colon \mathrm{input} \to \mathbb{R}$, решаемая на максимум или минимум. Для $\rho \geq 1$ алгоритм $g$ называется \emph{$\rho$-приближением}, если для любого входа $x$:
	\begin{align*}
		g(x) &\geq \frac{f^*(x)}{\rho} \quad \text{(для задачи максимизации),} \\
		g(x) &\leq f^*(x) \cdot \rho \quad \text{(для задачи минимизации),}
	\end{align*}
	где $f^*(x)$ --- оптимальное значение на входе $x$.
\end{Definition}

\begin{Remark}{}{}
	Величина $\rho$ называется \emph{коэффициентом приближения} (approximation ratio). Чем меньше $\rho$, тем лучше алгоритм. Значение $\rho = 1$ соответствует точному решению.
\end{Remark}

\subsection{Задача о вершинном покрытии}

\begin{Definition}{Вершинное покрытие}{}
	Дан граф $G = (V, E)$. \emph{Вершинное покрытие} --- это подмножество $C \subseteq V$ такое, что для каждого ребра $e \in E$ хотя бы один из его концов принадлежит $C$:
	\[
	\forall \{u,v\} \in E \;\; \exists\, w \in C \colon w \in \{u,v\}.
	\]
\end{Definition}

\textbf{Задача:} найти вершинное покрытие минимального размера $|C|$.

Задача о вершинном покрытии является NP-полной. Для неё неизвестно полиномиальных алгоритмов (в стандартной теории $\mathrm{P} \neq \mathrm{NP}$). Однако существует простой 2-приближённый алгоритм, основанный на жадном построении паросочетания.

\subsubsection*{Алгоритм жадного паросочетания (2-приближение)}

\begin{enumerate}
	\item Выбрать произвольное ребро $\{u,v\} \in E$.
	\item Добавить обе вершины в покрытие: $C \coloneqq C \cup \{u, v\}$.
	\item Удалить из графа $u$, $v$ и все инцидентные им рёбра: $G \coloneqq G \setminus \{u,v\}$.
	\item Если $E(G) \neq \varnothing$, перейти к шагу~1.
\end{enumerate}

На выходе $C$ является вершинным покрытием, поскольку каждое ребро было <<обработано>> на некотором шаге.

\begin{Statement}{Алгоритм жадного паросочетания даёт 2-приближение}{}
	Пусть $C^*$ --- минимальное вершинное покрытие. Тогда $|C| \leq 2\,|C^*|$.
\end{Statement}

\begin{proof}
	В ходе алгоритма строится набор рёбер $M$ (паросочетание): на каждом шаге выбирается ребро $\{u,v\}$, причём после удаления $u$ и $v$ ни одно из последующих рёбер не касается $u$ или $v$. Таким образом, $M$ --- паросочетание размера $|C|/2$.
	
	Так как $C^*$ является вершинным покрытием, каждое ребро из $M$ должно иметь хотя бы один конец в $C^*$. Поскольку рёбра паросочетания не делят вершин, различным рёбрам из $M$ соответствуют различные вершины из $C^*$. Следовательно,
	\[
	|C^*| \geq |M| = \frac{|C|}{2} \implies |C| \leq 2\,|C^*|.
	\]
\end{proof}

\subsection{Задача о взвешенном вершинном покрытии}

Дан граф $G=(V,E)$ и функция весов $w\colon V \to \mathbb{R}_{>0}$. Требуется найти вершинное покрытие $C\subseteq V$ минимального суммарного веса:
\[
\sum_{v \in C} w(v) \to \min.
\]

Алгоритм жадного паросочетания в этом случае \emph{не работает} с прежней гарантией, поскольку не учитывает веса вершин. Задача о взвешенном вершинном покрытии также NP-полная (эквивалентна по трудности невзвешенному случаю).

\subsubsection*{Сведение к целочисленному линейному программированию}

Введём переменные $x_v \in \{0,1\}$ для каждой вершины $v$: $x_v = 1$ означает, что вершина $v$ включена в покрытие. Задача записывается как:
\[
\begin{cases}
	\displaystyle\min_{x} \sum_{v \in V} w(v)\cdot x_v, \\[4pt]
	x_u + x_v \geq 1 \quad \forall\, \{u,v\} \in E, \\[2pt]
	0 \leq x_v \leq 1 \quad \forall\, v \in V, \\[2pt]
	x_v \in \{0,1\}.
\end{cases}
\]

\subsubsection*{2-приближённый алгоритм через LP-релаксацию}

Снимем условие целочисленности и решим непрерывную \emph{LP-релаксацию} (без ограничения $x_v \in \{0,1\}$):
\[
\begin{cases}
	\displaystyle\min \sum_{v \in V} w(v)\cdot x_v, \\[4pt]
	x_u + x_v \geq 1 \quad \forall\, \{u,v\} \in E, \\[2pt]
	0 \leq x_v \leq 1 \quad \forall\, v \in V.
\end{cases}
\]
LP-релаксация решается за полиномиальное время, и мы находим оптимальное дробное решение $\{x^*_v\}$.

Построим целочисленное решение \emph{округлением}:
\[
C = \bigl\{\, v \in V \mid x^*_v \geq \tfrac{1}{2} \,\bigr\}, \qquad
\bar{x}_v = \begin{cases} 1, & v \in C, \\ 0, & v \notin C. \end{cases}
\]

\begin{Theorem}{О 2-приближении через LP-релаксацию}{}
	Множество $C$, построенное описанным выше методом, является вершинным покрытием и 2-приближением задачи о взвешенном вершинном покрытии.
\end{Theorem}

\begin{proof}
	\emph{Корректность покрытия.} Для каждого ребра $\{u,v\} \in E$ из ограничения LP: $x^*_u + x^*_v \geq 1$. Следовательно, хотя бы одно из двух неравенств $x^*_u \geq \tfrac{1}{2}$ или $x^*_v \geq \tfrac{1}{2}$ выполнено, то есть хотя бы одна из вершин попадает в $C$. Значит, $C$ --- вершинное покрытие.
	
	\emph{Оценка качества.} Так как $\bar{x}_v = 1$ только при $x^*_v \geq \tfrac{1}{2}$, имеем $\bar{x}_v \leq 2 x^*_v$ для всех $v$. Поэтому:
	\[
	\sum_{v \in V} w(v)\,\bar{x}_v \;\leq\; \sum_{v \in V} w(v) \cdot 2 x^*_v \;=\; 2\sum_{v \in V} w(v)\, x^*_v \;\leq\; 2\sum_{v \in V} w(v)\, x^{**}_v,
	\]
	где $\{x^{**}_v\}$ --- оптимальное целочисленное решение (оно допустимо для LP, поэтому $\sum w(v) x^*_v \leq \sum w(v) x^{**}_v$).
	
	Следовательно, $\sum_{v \in C} w(v) \leq 2\sum_{v \in V} w(v)\,x^{**}_v$.
\end{proof}

\begin{Remark}{}{}
	Эта задача является ещё и более оптимальной: известно, что $(2-\varepsilon)$-приближение для вершинного покрытия крайне сложно или, по-видимому, невозможно найти (при стандартных предположениях о сложности).
\end{Remark}